\begin{table}[H]
  \centering
  \small
  \caption{Description textuelle du cas d'utilisation : «~Soumettre une demande de modification pour un plan~»}
  \label{tab:cu-soumettre-demande-modification}
  \PFETableSetup
  \PFETableRowColors
  \PFETableArrayStretch
  \PFETableTabularXBegin
  \begin{tabularx}{\PFETableBlockWidth}{@{}>{\raggedright\arraybackslash}p{3.2cm} >{\raggedright\arraybackslash}X@{}}
    \tableHeaderStyle
    \tableHeaderCell{Champ} & \tableHeaderCell{Contenu} \\

    \textbf{Titre} & Soumettre une demande de modification pour un plan \\

    \textbf{Acteurs} & Utilisateur Site \\

    \textbf{Objectif} &
      Permettre à l'utilisateur de formuler une demande de modification sur un plan verrouillé, afin de corriger ou ajuster des informations après validation. \\

    \textbf{Préconditions} &
      1. L'utilisateur est authentifié.\par
      2. Il est rattaché au site concerné.\par
      3. Le plan existe.\par
      4. Le plan est dans un état non modifiable. \\

    \textbf{Postconditions} &
      1. La demande de modification est enregistrée.\par
      2. Son statut est défini à «~En attente~».\par
      3. Une notification est envoyée à l'utilisateur Corporate.\par
      4. Une trace est enregistrée dans l'historique. \\

    \textbf{Scénario nominal} &
      1. L'utilisateur accède au plan concerné.\par
      2. Le système affiche les détails du plan.\par
      3. L'utilisateur sélectionne l'option de demande de modification.\par
      4. Le système affiche le formulaire de demande.\par
      5. L'utilisateur saisit les informations requises (raison, durée, ...).\par
      6. L'utilisateur joint une pièce justificative si nécessaire.\par
      7. Le système vérifie la complétude des données saisies.\par
      8. Le système enregistre la demande de modification.\par
      9. Le système déclenche une notification vers l'utilisateur Corporate.\par
      10. Le système enregistre l'action dans l'historique.\par
      11. La demande apparaît dans la liste des demandes en attente. \\

    \textbf{Scénario alternatif} &
      \textbf{1. Données incomplètes :}\par
      1.1 L'utilisateur soumet une demande avec des champs obligatoires manquants.\par
      1.2 Le système refuse la soumission.\par
      1.3 Un message d'erreur précise les informations à compléter. \\
  \end{tabularx}%
  \PFETableTabularXEnd
\end{table}
